\chapter*{Appendix II}

\addcontentsline{toc}{chapter}{Appendix II}

\def\shorttitle{Appendix II}

\section*{Glossary of common terms}

\textit{System}: Grouping of atoms, minerals, rocks and/or gases and waters under consideration within a single volume of space (Langmuir, 1997). The extent of a system may be defined as is convenient and can vary from a droplet of rain, a unit volume of aquifer to an ocean. An open system can exchange mass and energy with its surroundings. A closed system can only exchange energy.

\textit{Phase}: A region of space that is physically distinct, mechanically separable and homogeneous in its composition and properties (Bethke, 1996). A homogeneous system contains only one phase, a heterogeneous system contain more than one phase. Geochemical models always contain at least one aqueous phase. Additionally, gaseous or mineral phases may be considered. Rocks such as granites or feldspar do not constitute phases. They consist of different (mineral) phases.

\textit{Species}: Molecular entities that exist within a phase (Bethke, 1996). Examples are \co\ or \ox\ in a gas, \na\ or \su\ in an aqueous solution or a complex like \si.

\textit{Component}: Simple chemical units that can be combined to describe the chemical composition of the species or substances in a system (Langmuir, 1997). Components can be thought of as the ingredients in a recipe. Whereas species and phases exist as real entities, components are merely mathematical tools for describing the system�s composition (Bethke, 1996). Components may also occur as species and it�s very important to discern between the two. For example, bicarbonate (\bi) may be chosen as a component and is distributed among the species: \bi, \co, CO32- CaHCO3+, NaCO3-.
Components are termed \textit{master species} in PHREEQC.

\textit{Basis}: The set of components used in a geochemical model. The choice of a basis is arbitrarily and there is no single basis that describes a system (Bethke, 1996). Three rules must be followed in choosing the basis:

\begin{itemize}
\item Each species and phase in our model can be formed from a combination of components in the basis
\item The number of components in the basis is the minimum necessary to satisfy the first rule
\item It is impossible to write a balanced reaction to form a component in terms of other components. This means that the components must be linearly independent of one another.
\end{itemize}